% =============================================================================
% P3D -- Lições aprendidas (PT-BR)
% Formatação: ABNT (classe abnTeX2, modo article), A4, Latin Modern, 12 pt,
% margens 3/2/3/2 cm, espaçamento 1,5, recuo de parágrafo 1,25 cm.
% Figuras esperadas em out/ (compiladas a partir de figuras/):
%   Fig1_PT.pdf, Fig4_PT.pdf, Fig5_PT.pdf
% (as figuras 2 e 3 são tabelas coloridas escritas neste arquivo).
% ARQUIVO GERADO por ferramentas/gerar_tex.py a partir de ../MD. Para mudar o
% conteúdo, edite o Markdown e gere de novo, ou edite este .tex diretamente.
% =============================================================================
\documentclass[12pt,a4paper,article,oneside,english,brazil]{abntex2}

\usepackage[T1]{fontenc}
\usepackage[utf8]{inputenc}
\usepackage{lmodern}           % Latin Modern (fonte padrão do LaTeX, com T1)
\usepackage{amsmath}
\usepackage{graphicx}
\usepackage{xcolor}
\usepackage{array}
\usepackage{colortbl}
\usepackage{tabularx}
\usepackage{longtable}
\usepackage[expansion=false]{microtype}
\usepackage{url}
\Urlmuskip=0mu plus 1mu\relax

% Margens ABNT: superior 3 cm, inferior 2 cm, esquerda 3 cm, direita 2 cm
\setlrmarginsandblock{3cm}{2cm}{*}
\setulmarginsandblock{3cm}{2cm}{*}
\checkandfixthelayout

\setlength{\parindent}{1.25cm}
\setlength{\parskip}{0pt}
\setlength{\emergencystretch}{3em}
\setcounter{secnumdepth}{3}
\setcounter{tocdepth}{3}
\renewcommand{\contentsname}{SUMÁRIO}
\renewcommand{\listfigurename}{LISTA DE ILUSTRAÇÕES}
\renewcommand{\listtablename}{LISTA DE TABELAS}

% Cores das figuras-tabela (mapa de calor e matriz de princípios)
\definecolor{exLine}{RGB}{23,54,93}
\definecolor{bandaProativa}{RGB}{220,233,247}
\definecolor{bandaResp}{RGB}{200,240,236}
\definecolor{bandaValor}{RGB}{238,242,245}

% Tabelas: sem filetes verticais (padrão IBGE adotado pela ABNT)
\renewcommand{\arraystretch}{1.2}

% Fonte abaixo de tabelas e figuras (ABNT NBR 14724)
\newcommand{\fontetab}[1]{\par\vspace{3pt}{\small\noindent Fonte: #1\par}}

% Citações numéricas: o número de cada referência é definido uma vez (ordem
% alfabética da lista) e usado por \refnum{chave}, \cit{chave} e \refitem.
\newcommand{\defref}[2]{\expandafter\gdef\csname refnum@#1\endcsname{#2}}
\newcommand{\refnum}[1]{\csname refnum@#1\endcsname}
\newcommand{\cit}[1]{[\refnum{#1}]}
\newcommand{\refitem}[2]{\par\noindent\hangindent=1.2cm\hangafter=1\makebox[1.2cm][l]{[\refnum{#1}]}#2\par}


% Integrantes do grupo (preencher; a capa usa esta macro)
\newcommand{\integrantes}{[nomes dos integrantes do grupo]}

\begin{document}
\SingleSpacing
\thispagestyle{plain}

\begin{center}\textbf{RELATO DE LIÇÕES APRENDIDAS}\end{center}

\textbf{Grupo SIPA Privacidade 2}

\subsection*{O que aprendemos na síntese do PMBOK 8}

\textbf{1. Separar princípio, área de foco, domínio e processo.} Esses quatro elementos respondem a perguntas diferentes: o princípio orienta julgamento; a área de foco situa um tipo de atividade; o domínio reúne uma preocupação de gestão; o processo descreve uma prática possível. Para não sugerir uma cadeia rígida entre eles, usamos verbos distintos nas relações e matrizes separadas: 7 \ensuremath{\times} 5 para processos e áreas de foco, e 6 \ensuremath{\times} 7 para princípios e domínios.

\textbf{2. Registrar a origem de cada contagem.} O texto convertido da Tabela 2-1 do \textit{Guide} lista os processos de cada domínio sem indicar a área de foco. A área de \textit{Identify Stakeholders} vem da Figura 2-32 do \textit{Guide}, que o coloca em Iniciação. Assim, os totais por área são 2, 19, 8, 10 e 1, e os totais por domínio são 9, 6, 3, 4, 7, 5 e 6. Anotar a fonte de cada número tornou o quadro verificável.

\textbf{3. Fazer do \textit{tailoring} o eixo da síntese.} Uma lista dos 40 nomes, sozinha, seria pouco explicativa. O capítulo 3 do \textit{Guide} mostra por que e como ajustar ciclo de vida, processos e engajamento, e traz diagnósticos de situações comuns. Passamos a ler cada domínio pela pergunta: que decisão de adaptação este domínio exige no contexto do projeto?

\textbf{4. Distinguir o PMBOK 8 do nosso agrupamento analítico.} As quatro concentrações ajudam a apresentar os sete domínios, mas não são classificação oficial do PMI. Os temas A, S, I e P são acréscimos em análise e não entram na soma dos 40 processos. Tornar essa fronteira visível nas legendas e no texto evitou atribuição incorreta à fonte.

\textbf{5. Trabalhar com uma fonte de acesso restrito.} A edição integral do PMBOK 8 só se obtém mediante aquisição ou associação ao PMI. Por isso o relatório cita seções, tabelas e figuras, parafraseia em vez de reproduzir trechos e deixa claro que as traduções dos nomes dos processos são de trabalho.

\textbf{6. Completar o PMBOK 8 com literatura externa.} A edição não menciona gêmeo digital. Estudos recentes mostram que abordagens híbridas são comuns e que a escolha depende do contexto, o que reforça a ênfase do PMBOK 8 na adaptação. Para projetos de gêmeo digital, a estrutura de gestão do PMBOK 8 precisa ser complementada com referências técnicas específicas.

\textbf{7. Usar figuras para explicar relações, não para decorar.} O mapa de calor mostra distribuição, o ciclo representa a revisão de decisões e os cartões de concentração deixam claro o que é oficial e o que é organização analítica. Cada figura tem título, fonte, legenda e frase interpretativa; os números e rótulos continuam legíveis sem cor.

\subsection*{O que revisaríamos numa próxima versão}

Conferir a tradução dos nomes dos processos com o texto original em inglês, testar se alguém que não participou da redação entende a diferença entre os dois volumes e obter validação de especialistas em gestão de projetos. O tamanho final das figuras depende da diagramação.

\textbf{Limite deste trabalho:} é uma síntese documental. Não relata um projeto concluído, não valida uma extensão do PMBOK 8 e não mede resultados dos exemplos hipotéticos.

\end{document}
